\section{Dataset Description}

\subsection{Dataset Composition}

Our dataset comprises 50,000+ English-Turkish parallel segments collected from diverse software localization sources. Table \ref{tab:dataset_stats} presents the composition and characteristics.

\begin{table}[h]
\centering
\caption{Dataset Statistics}
\label{tab:dataset_stats}
\begin{tabular}{@{}lrrr@{}}
\toprule
\textbf{Category} & \textbf{Segments} & \textbf{Avg. Length} & \textbf{Source} \\ \midrule
Technical Docs & 25,000 & 120.5 & OPUS, GitHub \\
UI Strings & 15,000 & 35.2 & Mozilla, Common UI \\
Error Messages & 10,000 & 55.8 & Mixed \\ \midrule
\textbf{Total} & \textbf{50,000} & \textbf{85.3} & \textbf{Multiple} \\ \bottomrule
\end{tabular}
\end{table}

\subsection{Data Sources}

\subsubsection{OPUS Corpus}

The Open Parallel Corpus (OPUS) \cite{tiedemann2012parallel} provides extensive multilingual parallel texts from open-source software projects. We extracted:

\begin{itemize}
    \item GNOME desktop environment documentation (10,000 segments)
    \item KDE system messages and manuals (8,000 segments)
    \item Mozilla Firefox localization files (6,000 segments)
\end{itemize}

\subsubsection{GitHub Open Source Projects}

We collected i18n files from popular Turkish software projects and bilingual README files, yielding 8,000 segments of authentic technical and UI content.

\subsubsection{MaCoCu Corpus}

From the MaCoCu Turkish-English corpus \cite{macocu2021} containing 1.6M parallel sentences, we filtered 10,000 segments with technical terminology and quality scores > 0.8.

\subsection{Dataset Characteristics}

\subsubsection{Text Length Distribution}

Figure \ref{fig:length_dist} shows the text length distribution across categories. Technical documentation exhibits longer segments (mean: 120.5 characters) compared to UI strings (mean: 35.2 characters), reflecting the different nature of these text types.

\subsubsection{Placeholder Analysis}

Approximately 25\% of UI strings contain placeholders (e.g., \texttt{\{username\}}, \texttt{\%d items}), which must be preserved during translation. Technical documentation contains fewer placeholders (8\%) but includes code snippets and technical notation.

\subsubsection{Terminology Density}

We identified 500+ unique technical terms (e.g., "function", "variable", "database", "authentication") appearing across the dataset. Technical documentation shows higher terminology density (15\% of tokens) compared to UI strings (5\%).

\subsection{Data Quality}

All source data consists of professional human translations from established open-source projects, ensuring high-quality reference translations for evaluation. We manually inspected random samples (n=200) and found 98\% accuracy in the reference translations.

\subsection{Train/Test Split}

We randomly split the dataset into:
\begin{itemize}
    \item \textbf{Training set}: 40,000 segments (80\%) - for future fine-tuning experiments
    \item \textbf{Test set}: 10,000 segments (20\%) - for evaluation in this study
\end{itemize}

The split maintains category proportions to ensure representative evaluation across text types.

\subsection{Dataset Availability}

The complete dataset, preprocessing scripts, and data statistics are publicly available at [GitHub URL] under [License]. This enables reproduction of our experiments and facilitates future research on Turkish software localization.
